% 19 November 2026. Prepared for Edward Baker with OpenAI Codex assistance.
% Lead recorded gpt-6-astra/max. Inherited results retain their source scope.
\section{The source, the time direction and the question}
\label{sec:source}

\subsection{A distinct use of the existing hierarchy}

The present investigation begins with a concrete question. Suppose a
source theorem certifies separated zeros at a positive heat time. Can
arithmetic information about the error in their evolution preserve that
separation at an earlier time where the initial localization theorem
no longer applies? The question concerns the usefulness of a quantitative
estimate along the intervening path. A formula for the velocity at one
time, or a more accurate description of an already protected region,
is not yet that estimate.

Our starting framework is the exterior-moment manuscript
\cite{VCDCMManuscript}. That work develops exact pair dynamics,
compatible common-time approximation and direct contour enclosures.
Here we select and rederive the background needed to connect those
identities to the arithmetic approximation of Polymath
\cite{Polymath2019}. The motivating discussion
\cite{VCDCMDiscussion08,VCDCMDiscussion09,VCDCMDiscussion10,
VCDCMDiscussion11,VCDCMDiscussion12} supplies three guiding distinctions:
actual gaps versus a reference density, a joint drift versus separate
moment errors, and approximation error versus the final inequality
which protects a zero configuration. The new program is organized around
these distinctions rather than around another representation alone.

\subsection{The literal theta field}

For a source coordinate $v\ge0$ put
\begin{equation}
 \Phi(v)=\sum_{n=1}^{\infty}
 \left(2\pi^2n^4e^{9v}-3\pi n^2e^{5v}\right)
 e^{-\pi n^2e^{4v}},
 \qquad \Phi(-v)=\Phi(v).
 \label{src:phi}
\end{equation}
For real physical time $t$, and then by analytic continuation, define
\begin{equation}
 H(t,z)=\int_0^{\infty}e^{tv^2}\Phi(v)\cos(zv)\dd v
       =\frac12\int_{\R}e^{tv^2}\Phi(v)e^{\ii zv}\dd v.
 \label{src:H}
\end{equation}
The source and observation kernel remain fixed throughout this paper.
With the usual completion
\[
 \xi(s)=\tfrac12s(s-1)\pi^{-s/2}\Gamma(s/2)\zeta(s),
\]
the theta Mellin representation gives
\begin{equation}
 H(0,z)=\frac18\xi\!\left(\frac{1+\ii z}{2}\right).
 \label{src:xi}
\end{equation}
This normalization is the one used in
\cite{Polymath2019,VCDCMManuscript}. In particular Newman height $x$
is twice the corresponding absolute zeta height, not the zeta height
itself. The functional equation and evenness make either sign of $\ii z$
in \eqref{src:xi} equivalent.

\begin{proposition}[Analyticity and heat law]
\label{src:heat}
The field $H$ is jointly entire in $(t,z)$, even in $z$, and real for
real $t,z$. It obeys
\begin{equation}
 H_t=-H_{zz}.
 \label{src:heat-law}
\end{equation}
For every compact real time interval $J$, $H(t,0)$ has a strictly
positive minimum and the spatial zero divisor satisfies uniformly
\begin{equation}
 n_t(R)=O(R\log(2+R)),\qquad
 \sum_{|\rho|>R}|\rho|^{-2}
       =O\!\left(\frac{\log(2+R)}R\right).
 \label{src:count}
\end{equation}
Zeros in the sum are counted with multiplicity.
\end{proposition}

\begin{proof}
Each term of \eqref{src:phi} is positive on $v\ge0$, since
$2\pi n^2e^{4v}-3>0$. The defining series and each fixed derivative
have a bound by a polynomial in $e^v$ times
$\exp(-c e^{4v})$ for some $c>0$. On compact $(t,z)$ sets this
majorizes the integrand and its derivatives after multiplication by
$\exp(\Re(t)v^2+|\Im z|v)$. Differentiation under the integral
proves the analytic assertions and \eqref{src:heat-law}. Positivity,
continuity and compactness give the lower bound at zero.

To control spatial growth, reserve half the double-exponential decay
for integrability. The supremum of $rv-c e^{4v}$ is
$O(r\log(2+r))$; the time factor contributes only a lower-order
$O((\log(2+r))^2)$ at that maximizing scale, uniformly for $t\in J$.
Thus $\log\max_{|z|\le r}|H(t,z)|=O(1+r\log(2+r))$.
Jensen's formula, using the positive lower bound at zero and the
maximum on $|z|=2R$, proves the count. Stieltjes summation by parts
bounds the inverse-square tail by a constant times
$\int_R^\infty \log(2+s)s^{-2}\dd s$.
\end{proof}

The growth is of order at most one. Evenness therefore pairs the
canonical factors for opposite zeros, removing their linear
exponentials. Proposition~\ref{src:heat} supplies the complete tail
needed for the rational traces below. It does not locate zeros,
separate them, or impose a positive sign on an arbitrary exterior
interaction.

\subsection{Backward time is a forward heat equation}

Fix a reference time $t_*$ and put $u=t_*-t$. Then
\begin{equation}
 \partial_u H=H_{zz}.
 \label{src:backward}
\end{equation}
The words ``backward continuation'' refer to decreasing physical Newman
time. They do not mean reversing the usual positive-diffusion equation
in the real spatial coordinate. For $u\ge0$ the exact field satisfies
\begin{equation}
 H(t_*-u,z)=\frac1{\sqrt\pi}\int_{\R}
 e^{-r^2}H(t_*,z+2\sqrt u\,r)\dd r.
 \label{src:gaussian}
\end{equation}
To verify this formula, insert \eqref{src:H}, interchange integrals
using the absolute source bound, and evaluate the Gaussian transform
$\pi^{-1/2}\int e^{-r^2}e^{2\ii\sqrt u\,rv}\dd r=e^{-uv^2}$.
This is a whole-line identity. Restricting the initial data to a small
interval without paying the complement does not preserve it.

The de Bruijn--Newman constant $\Lambda$ is characterized by the
all-real-zero regime $t\ge\Lambda$; the Riemann hypothesis is equivalent
to $\Lambda\le0$. The result $\Lambda\ge0$ of Rodgers and Tao
\cite{RodgersTao2020} supplies context, not an estimate used in the local
arguments. Our aim is a useful source-specific continuation estimate;
no assertion about the value of $\Lambda$ follows from the setup.

\subsection{Notation and evidence}

\begin{center}
\begin{tabular}{>{\raggedright\arraybackslash}p{0.18\textwidth}
                >{\raggedright\arraybackslash}p{0.73\textwidth}}
\toprule
Symbol & Meaning \\
\midrule
$t,u$ & Physical Newman time and reversed time $u=t_*-t$. \\
$a,b,c,g,D$ & Pair endpoints, center, actual gap $g=b-a$, and discriminant $D=g^2/4$. \\
$Q_j,M$ & Exterior denominator $(\rho_j-c)^2-D$ and complete interaction $\sum_j Q_j^{-1}$. \\
$d,\ell$ & Height-dependent reference spacing and a separately frozen normalization scale. Neither is an actual gap by definition. \\
$p,q$ & Nonnegative integer split-pole exponents unless explicitly extended. A single complex exponent later acts on the combined denominator. \\
$S_n,T_n$ & Even and odd complete exterior moments. \\
$\mathcal E,V,K$ & Gap drift, its half-defect, and the dimensionless joint forcing. \\
$P,E,\mathcal R$ & Arithmetic reference field, error $E=H-P$, and heat defect $\mathcal R=P_t+P_{zz}$. \\
\bottomrule
\end{tabular}
\end{center}

A local identity is proved under explicit product and chart assumptions.
A continuation theorem requires inequalities along a stopped time path.
A canonical acquisition proves those inequalities for the unchanged theta
source. An improved global bound additionally requires complete coverage.
The manuscript keeps these four levels distinct. Its derivations are
internal LLM-assisted work built on the cited sources; publication or
independent specialist review is not asserted.
